\documentclass[11pt,a4paper]{article}
\usepackage[utf8]{inputenc}
\usepackage[margin=1in]{geometry}
\usepackage{amsmath,amssymb,amsthm}
\usepackage{listings}
\usepackage{xcolor}
\usepackage{hyperref}

\newtheorem{theorem}{Theorem}
\newtheorem{lemma}[theorem]{Lemma}

\lstset{
  language=C++,
  basicstyle=\ttfamily\small,
  keywordstyle=\color{blue}\bfseries,
  commentstyle=\color{green!60!black},
  stringstyle=\color{red!70!black},
  numbers=left,
  numberstyle=\tiny\color{gray},
  breaklines=true,
  frame=single,
  tabsize=4,
  showstringspaces=false
}

\title{IOI 2010 -- Traffic}
\author{Solution}
\date{}

\begin{document}
\maketitle

\section{Problem Summary}
Given a tree of $N$ cities where city~$i$ has population $p_i$, find a city $r$ such that when the tree is rooted at $r$, the maximum subtree population among $r$'s children is minimized. This node is called the \emph{traffic center}.

\section{Solution}

\subsection{Key Observation}
Root the tree at node~$0$. Let $\mathrm{sub}(u)$ denote the total population in $u$'s subtree, and let $T = \sum_i p_i$ be the total population.

If the tree is re-rooted at node $v$, the ``directions'' from $v$ are:
\begin{itemize}
    \item Each child $c$ of $v$ (in the original rooting) contributes $\mathrm{sub}(c)$ population.
    \item The ``parent direction'' contributes $T - \mathrm{sub}(v)$ population.
\end{itemize}

The maximum directional load at $v$ is therefore:
\[
f(v) = \max\!\bigl(T - \mathrm{sub}(v),\; \max_{c \in \mathrm{children}(v)} \mathrm{sub}(c)\bigr).
\]

The traffic center is $\arg\min_v f(v)$.

\subsection{Algorithm}
\begin{enumerate}
    \item Root at node~$0$. Compute $\mathrm{sub}(u)$ for all $u$ via reverse BFS.
    \item For each node $u$, compute $\max_{c} \mathrm{sub}(c)$ (the maximum child subtree population).
    \item For each node $v$, compute $f(v)$ and track the minimizer.
\end{enumerate}

\begin{theorem}
This computes the traffic center in $O(N)$ time and space.
\end{theorem}
\begin{proof}
Steps 1--3 each take a single pass over the tree. The formula for $f(v)$ correctly captures the maximum directional load at $v$ because, in the re-rooted tree, every direction from $v$ either corresponds to a child subtree or the complementary set $T - \mathrm{sub}(v)$.
\end{proof}

\subsection{Complexity}
\begin{itemize}
    \item \textbf{Time:} $O(N)$.
    \item \textbf{Space:} $O(N)$.
\end{itemize}

\section{Implementation}

\begin{lstlisting}
#include <bits/stdc++.h>
using namespace std;

int main() {
    ios::sync_with_stdio(false);
    cin.tie(nullptr);

    int N;
    cin >> N;

    vector<long long> pop(N);
    for (int i = 0; i < N; i++) cin >> pop[i];

    vector<vector<int>> adj(N);
    for (int i = 0; i < N - 1; i++) {
        int u, v;
        cin >> u >> v;
        adj[u].push_back(v);
        adj[v].push_back(u);
    }

    // Root at 0, BFS to get order and parents
    vector<long long> sub(N, 0);
    vector<int> parent(N, -1), order;
    vector<bool> visited(N, false);
    queue<int> q;
    q.push(0);
    visited[0] = true;
    while (!q.empty()) {
        int u = q.front(); q.pop();
        order.push_back(u);
        for (int v : adj[u])
            if (!visited[v]) {
                visited[v] = true;
                parent[v] = u;
                q.push(v);
            }
    }

    // Compute subtree populations in reverse BFS order
    for (int i = 0; i < N; i++) sub[i] = pop[i];
    for (int i = N - 1; i >= 1; i--)
        sub[parent[order[i]]] += sub[order[i]];

    long long total = sub[0];

    // Compute max child subtree population for each node
    vector<long long> maxChild(N, 0);
    for (int i = 1; i < N; i++)
        maxChild[parent[order[i]]] = max(maxChild[parent[order[i]]],
                                         sub[order[i]]);

    // Find traffic center
    long long bestLoad = LLONG_MAX;
    int bestNode = 0;
    for (int u = 0; u < N; u++) {
        long long load = max(total - sub[u], maxChild[u]);
        if (load < bestLoad) {
            bestLoad = load;
            bestNode = u;
        }
    }

    cout << bestNode << "\n";
    return 0;
}
\end{lstlisting}

\end{document}
